\chapter{State of the Problem}

\section{Current Approaches in Canine Gait Assessment}

Existing methods for gait analysis in animals can be broadly divided into
laboratory-based systems, wearable systems, and vision-based approaches.

Laboratory-based systems, such as force plates and pressure walkways, provide
accurate measurements of ground reaction forces and spatiotemporal
parameters~\cite{McLaughlin2001,Oosterlinck2011,Fahie2018}.
However, these systems require controlled environments, specialised facilities,
and experienced operators, making them unsuitable for continuous monitoring in
real-world conditions or for use in veterinary clinics without dedicated
infrastructure.

Wearable systems, typically based on inertial sensors, enable data collection
during natural movement~\cite{Jenkins2018,Ladha2017,Zhang2022,Hayati2019}.
These systems are more flexible and can be deployed in real-world conditions,
but are generally limited to kinematic measurements and do not capture the
kinetic information necessary to assess limb loading and prosthetic force
transmission.

Vision-based approaches rely on video analysis to extract motion parameters.
While useful for qualitative assessment, these methods require complex setups
and are sensitive to environmental conditions, limiting their applicability
outside controlled laboratory environments.

\section{Instrumented Systems for Locomotion Analysis}

Recent work has explored the integration of inertial and force sensing to
provide a more complete description of locomotion, capturing both kinematic and
kinetic data~\cite{Zhang2022}.
In veterinary applications, such systems have been used to detect lameness,
characterise gait patterns, and monitor rehabilitation
processes~\cite{Fahie2018,Blattler2024}.
However, their integration into prosthetic devices remains limited, and their
use is generally constrained to short experimental sessions rather than
continuous monitoring.

\section{Comparative Analysis of Existing Approaches}

Table~\ref{tab:state_of_art} provides a structured comparison of representative
approaches from the literature along dimensions relevant to the proposed system.
These dimensions reflect the key challenges identified in veterinary prosthetic
evaluation: real-world deployability, kinetic sensing capability,
single-limb applicability, and integration with a prosthetic device.

\begin{table}[ht]
  \centering
  \caption{Comparison of existing approaches for canine locomotion assessment.}
  \label{tab:state_of_art}
  \small
  \begin{tabular}{p{3.0cm} p{2.0cm} c c c c}
    \toprule
    \textbf{Reference}
      & \textbf{Sensing}
      & \textbf{Kinetics}
      & \textbf{Wearable}
      & \textbf{Single limb}
      & \textbf{Prosthetic} \\
    \midrule
    McLaughlin~\cite{McLaughlin2001}
      & Force plate     & \checkmark & --         & --         & -- \\
    Oosterlinck~\cite{Oosterlinck2011}
      & Pressure mat    & \checkmark & --         & --         & -- \\
    Fahie~\cite{Fahie2018}
      & Pressure walkway & \checkmark & --        & --         & -- \\
    Jenkins~\cite{Jenkins2018}
      & IMU (1)         & --         & \checkmark & \checkmark & -- \\
    Ladha~\cite{Ladha2017}
      & IMU (1)         & --         & \checkmark & \checkmark & -- \\
    Zhang~\cite{Zhang2022}
      & IMU (4)         & --         & \checkmark & --         & -- \\
    Hayati~\cite{Hayati2019}
      & IMU (1)         & --         & \checkmark & \checkmark & -- \\
    Bl\"{a}ttler~\cite{Blattler2024}
      & IMU (4)         & --         & \checkmark & --         & -- \\
    \midrule
    \textbf{This work}
      & IMU + load cell & \checkmark & \checkmark & \checkmark & \checkmark \\
    \bottomrule
  \end{tabular}
\end{table}

\noindent
The table shows that no existing approach combines all five properties.
Laboratory-based systems (force plates, pressure mats) provide kinetic
measurements but are not wearable and cannot be deployed in real-world
conditions. Wearable IMU-based systems are portable and suitable for
single-limb measurement, but do not capture force data.
Critically, no existing work integrates sensing directly into a prosthetic
device to enable continuous and autonomous monitoring.

\section{Limitations of Existing Solutions}

Despite significant advances, current solutions present several limitations
that are particularly relevant in the context of prosthetic evaluation.

Laboratory-based systems lack portability and cannot be used in everyday
environments or across extended monitoring periods.
Wearable systems distributed across limbs increase complexity and reduce
practicality for long-term use~\cite{Zhang2022}.
Most systems are not designed for prosthetic integration and therefore
cannot assess force transmission or load acceptance in the prosthetic device.

A fundamental methodological limitation is the reliance on inter-limb
symmetry metrics, which require simultaneous measurement of multiple
limbs~\cite{Oosterlinck2011,Fahie2018}.
This assumption is not valid in prosthetic applications where only the
prosthetic limb is instrumented.
Furthermore, wearable systems designed for healthy dogs typically require
sensors to be attached directly to the animal's body, increasing complexity
and stress during instrumentation~\cite{Zhang2022,Blattler2024}.

\section{Research Gap}

The analysis of existing approaches reveals a clear gap: there is no
portable, integrated system capable of acquiring continuous kinematic and
kinetic data directly from a prosthetic limb in real-world conditions.

This work addresses this gap by embedding inertial and force sensing directly
into the prosthetic structure, enabling continuous and objective assessment of
prosthetic adaptation using metrics derived exclusively from the prosthetic limb.
